
\section{引言}
\subsection{问题背景}

零售业是国民经济的重要组成部分，直接关系到消费市场的活跃度与经济增长质量。随着电子商务的迅猛发展和消费升级趋势的深化，传统零售企业面临着日益复杂的市场环境和激烈的竞争压力。

然而，当前许多零售企业在经营决策中仍高度依赖经验判断，缺乏数据驱动的科学预测工具。具体表现为：难以准确识别销售额的关键影响因素，导致资源配置效率低下；缺乏高精度的销售预测模型，造成库存积压与缺货并存；在面对市场波动时缺乏柔性应对策略，损失潜在的营收机会。

因此，亟需建立一套基于数据分析的零售销售额预测与优化体系，帮助企业精准把握销售规律、科学制定经营策略。本文基于某零售企业的日常经营数据，围绕销售额影响因素识别、预测模型构建、模型评估优化和经营策略建议四个核心问题展开系统研究。

\subsection{问题重述}

\textbf{问题1：}基于给定的零售企业经营数据（包含客流量、促销力度、客单价、线上占比等维度），识别影响销售额的关键因素，并量化各因素的影响程度。

\textbf{问题2：}建立数学模型对未来一段时间的销售额进行预测，要求模型具有较高的预测精度和良好的泛化能力。

\textbf{问题3：}对预测模型进行全面评估，包括交叉验证误差分析和关键参数的灵敏度分析，验证模型的可靠性和稳健性。

\textbf{问题4：}基于预测结果和敏感性分析，从客流管理、促销策略、渠道优化等方面为零售企业提供具体的经营策略建议。
